\documentclass[conference]{IEEEtran}
\IEEEoverridecommandlockouts

\usepackage{cite}
\usepackage{amsmath,amssymb}
\usepackage{graphicx}
\usepackage{booktabs}
\usepackage{url}
\usepackage[hidelinks]{hyperref}

\begin{document}

\title{Real-Time Sleep Staging Costs Less Than Reported:\\A Two-Dataset Evaluation-Protocol Study}

\author{\IEEEauthorblockN{Jahnvi R, Swapnil S, Shreeya S Methuku, Shreshtha Dixit}
\IEEEauthorblockA{Department of Computer Science (AIML) and Engineering,
PES University, Bengaluru, India}}

\maketitle

\begin{abstract}
A sleep-staging model on a wearable may use only past signal, and that constraint is assumed to
cost accuracy. We measure the cost with a matched ablation, one architecture trained with and
without access to future signal on identical folds and seeds, and repeat it on two datasets: 78
Sleep-EDF subjects and the 25 subjects of DOD-H. Evaluated the way this literature normally
evaluates, by tiling a recording with non-overlapping context windows, causality costs $0.025$
and $0.033$ Cohen's $\kappa$ and appears on both datasets to fall on REM. Most of that is
produced by the evaluation protocol. Tiling scores predictions at every position of the context
buffer, including positions with almost no history, which a streaming system passes through once
at start-up; the causal model is crippled there and the non-causal model is not. Moving the same
weights to a streaming protocol lifts the causal arm by $+0.0350$ and $+0.0345\,\kappa$ and the
non-causal arm by $+0.0036$ and $+0.0052$, an excess to the causal arm of $+0.0314$ and
$+0.0292$ that is positive in all 30 paired folds, with corrected $t(14) = 12.69$ and $6.48$. Once
that artifact is removed the residual cost is $+0.0065$ and $-0.0037\,\kappa$, neither resolvable,
and the REM concentration falls from $0.060$ F1 to $0.002$ and $0.023$. The model is strictly
causal from raw sample to logit, verified bit-for-bit, emits one label per second in 30{,}757
parameters at $0.026$~ms per second of EEG, and reports stage changes a median 19 seconds sooner
than an epoch-level system in 94.8\% and 95.6\% of transitions. For a device that must report as
it goes, the real-time constraint is close to free; the reported cost was mostly a measurement
artifact.
\end{abstract}

\begin{IEEEkeywords}
sleep staging, EEG, causal models, real-time inference, wearable devices, evaluation methodology
\end{IEEEkeywords}

\section{Introduction}

Clinical polysomnography is scored in 30-second epochs~\cite{aasm}, and most automated staging
models mirror that unit: they classify one epoch at a time and are often free to use signal from
later in the recording, or the whole night, to label an earlier epoch~\cite{deepsleepnet,
attnsleep}. That is appropriate for offline analysis of a completed study. It does not describe a
wearable device, which must report a stage as it happens if it is to drive an alert, a
closed-loop stimulus, or a live display.

Causality is therefore a real design constraint, and the natural question for anyone building
such a device is what it costs. That question is rarely answered cleanly, because published
causal and non-causal models differ in capacity, data, preprocessing and optimization all at
once. Output resolution is not the open problem: U-Sleep~\cite{usleep} and its predecessor
U-Time~\cite{utime} have emitted up to 7{,}680 stages per minute since 2021, using a
bidirectional model that reads several minutes of future signal, and the authors show that
high-frequency hypnograms carry clinical information that epoch-rate ones do not. What has not
been measured is what that resolution costs when the model may not look forward at all.

Causal staging is itself an active direction: Bresch et al.~\cite{bresch} apply recurrent
networks to real-time single-channel staging, CausalAttenNet~\cite{causalattennet} combines
causal convolution with long-term attention, and Micro SleepNet~\cite{microsleepnet} targets
mobile and wearable hardware. We claim no priority over these on the architecture. To our
knowledge none of them reports a matched ablation in which the identical network is trained with
and without access to future signal on identical folds and seeds, and none reports how the answer
moves with the evaluation protocol, which is what this paper is about.

We answer the question with such an ablation, run end to end on two datasets, and the useful
result is not the first number it produces but the reason that number is wrong.

Our contribution is one finding with three consequences. The finding is that batch-tiled
evaluation systematically penalizes causal models, measured as a within-arm quantity that does
not require the two arms to be comparable, and replicated on two corpora differing in electrode
montage, sampling rate, scoring team, class balance and subject count. The consequences are that
the reported causality penalty is largely an artifact of that protocol; that its apparent
concentration in REM is the same artifact; and that a corpus too small to resolve the between-arm
ablation can still resolve this effect decisively, which is what makes the measurement worth
adopting at the sample sizes this field actually works at.

\section{Methods}

\subsection{Datasets}

\textbf{Sleep-EDF-78}~\cite{physionet,kemp}: 78 subjects, 153 recordings, single channel Fpz-Cz at
100~Hz. Cassette recordings run about 20 hours with Wake at 68\%, so each night is trimmed to 30
minutes of Wake either side of its sleep period before a subject's recordings are concatenated,
leaving 1{,}629 hours with Wake at roughly 15\%.

\textbf{DOD-H}~\cite{dod}, the healthy-volunteer arm of the Dreem Open Datasets: 25 subjects,
205.5 hours, scored by five technicians, of which we use the released consensus. We use channel
F3-M2, one of only four frontal derivations present in both DOD-H and its clinical companion
DOD-O, resampled from 250 to 100~Hz with a polyphase anti-aliasing filter at the exact $2/5$ ratio
so the network's sample-defined kernels need no retuning. No Wake trimming is applied: DOD-H
records are 5--10 hours with Wake already at 12.3\%, and the Sleep-EDF setting would delete real
sleep-adjacent Wake rather than daytime recording. The class balance differs substantially from
trimmed Sleep-EDF (W 12.3\%, N1 6.1\%, N2 48.2\%, N3 14.2\%, REM 19.2\%), which is part of why it
is a useful second test.

Both datasets carry 30-second AASM labels~\cite{aasm}; each epoch's label is replicated across its
30 one-second training targets, so the network is supervised at 1~Hz even though the ground truth
is coarser. Per-second accuracy should be read as an approximation for that reason, and $\kappa$,
or accuracy re-aggregated to 30~s by majority vote, is the more comparable figure. Both use
subject-wise 5-fold cross-validation with no subject crossing the train/test boundary, verified by
unit test. Nothing about the architecture, the recipe or the two protocols changes between
datasets: the focal-loss class weights are recomputed from each training fold's own counts, so the
differing balances are handled automatically.

\subsection{Model and the Causality Ablation}

Two causal convolution branches read the raw signal in parallel, with 50-sample (0.5~s) and
400-sample (4~s) kernels targeting short events and slow waves respectively; concatenation gives 32
features per sample and left-aligned max-pooling reduces 100~Hz to one feature vector per second.
Three residual blocks of causal dilated convolutions~\cite{tcn} (dilations 1, 2, 4; 32 channels)
give a 29-second causal receptive field on that stream, and a four-head self-attention
layer~\cite{attention} with a causal mask extends effective history to the full 120-second context
window. A linear layer emits five class logits per second. Total 30{,}757 parameters.

Causality is enforced in exactly two places: convolutions pad only the left by $(K-1)d$, and the
attention mask sets $M_{ij} = -\infty$ for $j > i$. The non-causal arm removes both, padding
symmetrically and applying no mask, and changes nothing else, so the two arms are identical in
capacity and differ only in access to future signal.

Training uses focal loss~\cite{focal} ($\gamma = 2$, per-fold class weights), learning rate 0.001,
weight decay $10^{-4}$, batch size 64, and early stopping on validation $\kappa$ with patience 8.
The held-out fold is used for neither updates, stopping, nor checkpoint selection. Each arm is
trained under three seeds $\times$ five folds per dataset: 60 models in total.

\subsection{Causality Above the Network}

A network can be causal in every layer and still sit downstream of a preprocessing step that reads
the future. Per-epoch $z$-scoring, close to universal in this literature, normalizes a sample near
the start of an epoch using statistics computed over all 30 seconds of it. We replace it with a
trailing statistic computable one sample at a time,
\begin{equation}
z_t = \frac{x_t - \mu_t}{\sigma_t + \epsilon}, \qquad
\mu_t = \frac{1}{w}\sum_{i=t-w+1}^{t} x_i
\end{equation}
with $w = 30$~s and a reset at the start of each recording, so the opening seconds run on a partial
window exactly as a device powering on would. The offline array implementation and the
sample-at-a-time filter agree to $5.2\times10^{-14}$ over 24{,}000 samples.

Closing the leak improves both arms, from $\kappa = 0.6406$ to $0.6649$ causal and $0.6693$ to
$0.6898$ non-causal under tiled evaluation, so it was never a trade: the step that read the future
was costing accuracy rather than buying it. Every result below is produced under (1).

The full path from raw sample to logit is verified causal by perturbation test, and a streaming
front end reproduces the offline labels exactly over a 79{,}440-second held-out recording. Because
the two corpora are recorded in different amplitude units, we also verify that the trailing
normalizer removes a $1000\times$ change in input gain to within $5\times10^{-7}$, which is what
makes the pipeline portable between them without a unit conversion.

\subsection{The Two Evaluation Protocols}
\label{sec:protocol}

\textbf{Tiled.} The recording is covered by non-overlapping context windows and every position of
every window contributes one scored prediction. This is the default in every implementation we
have examined, including our own earlier work. Each window begins from an empty history.

\textbf{Streaming.} The model recomputes every 30~s against a full buffer and only the freshest 30
predictions of each window are kept, so after the opening window every scored second carries at
least 90 seconds of genuine history. This is what a deployed system does.

Every held-out second is scored exactly once under each protocol, from the same trained weights.
No model is retrained between protocols, which is what makes the comparison in
Section~\ref{sec:asym} a within-model one.

\subsection{Statistics}

Fifteen paired measurements per dataset. Because repeats reuse subjects across seeds, the primary
test is the Nadeau--Bengio corrected resampled $t$-test~\cite{nadeau}, which inflates the variance
estimate by $1/n + 1/(k-1)$, here $1/15 + 1/4$. Bootstrap intervals resample whole subjects rather
than seconds, since seconds within a recording are heavily autocorrelated. Pairing is exact: within
a seed the two arms use byte-identical subject lists, verified against the archived fold reports.

\section{Results}

\subsection{Headline Performance}

Table~\ref{tab:headline} gives the causal arm's absolute performance on both datasets under both
protocols. The streaming figures are the ones a deployed system would see: $\kappa = 0.700$ and
77.7\% per-second accuracy on Sleep-EDF-78, $\kappa = 0.667$ and 75.7\% on DOD-H.

\begin{table}[t]
\caption{Causal Arm, 15 Seed--Fold Cells Pooled Per Column}
\label{tab:headline}
\centering
\footnotesize
\setlength{\tabcolsep}{4pt}
\begin{tabular}{lcc|cc}
\toprule
 & \multicolumn{2}{c|}{Sleep-EDF-78} & \multicolumn{2}{c}{DOD-H} \\
Metric & Tiled & \textbf{Stream.} & Tiled & \textbf{Stream.} \\
\midrule
Accuracy (1~Hz)   & 0.7505 & \textbf{0.7772} & 0.7319 & \textbf{0.7574} \\
Cohen's $\kappa$  & 0.6649 & \textbf{0.6998} & 0.6329 & \textbf{0.6674} \\
Macro-F1          & 0.6894 & \textbf{0.7230} & 0.6765 & \textbf{0.7038} \\
N1 F1             & 0.4212 & \textbf{0.4629} & 0.3515 & \textbf{0.3938} \\
Accuracy at 30~s  & 0.7657 & \textbf{0.7823} & 0.7526 & \textbf{0.7659} \\
$\kappa$ at 30~s  & 0.6847 & \textbf{0.7065} & 0.6603 & \textbf{0.6788} \\
Fold $\kappa$ s.d.& 0.0257 & 0.0277 & 0.0982 & 0.1063 \\
\midrule
Subjects / hours & \multicolumn{2}{c|}{78 / 1{,}629} & \multicolumn{2}{c}{25 / 205.5} \\
\bottomrule
\end{tabular}
\end{table}

The fold-to-fold standard deviation on DOD-H is four times that on Sleep-EDF-78, a direct
consequence of holding out five subjects per fold rather than sixteen. That figure governs what can
and cannot be resolved on that dataset, and Section~\ref{sec:estimand} returns to it.

\subsection{The Causality Cost Depends on the Protocol}

Table~\ref{tab:protocol} reports the ablation under each protocol. On Sleep-EDF-78 the tiled
penalty is resolved and the streaming one is not, which is the shape the argument predicts. On
DOD-H neither is resolved: the corrected test returns $p = 0.22$ on a point estimate larger than
Sleep-EDF-78's, because five held-out subjects per fold is not enough. Read alone, this table would
rest on one dataset and be merely consistent with the other.

\begin{table}[t]
\caption{Cost of Causality Under Each Protocol. 15 Paired Folds Per Row, Nadeau--Bengio Corrected}
\label{tab:protocol}
\centering
\footnotesize
\setlength{\tabcolsep}{4pt}
\begin{tabular}{llccccc}
\toprule
Dataset & Protocol & Causal & Non-c. & $\Delta$ & Loses & $p$ \\
\midrule
Sleep-EDF-78 & Tiled     & 0.6649 & 0.6898 & $-0.0249$ & 14/15 & 0.005 \\
Sleep-EDF-78 & Streaming & \textbf{0.6998} & 0.6934 & $\mathbf{+0.0065}$ & 7/15 & 0.47 \\
\midrule
DOD-H        & Tiled     & 0.6329 & 0.6659 & $-0.0329$ & 10/15 & 0.22 \\
DOD-H        & Streaming & \textbf{0.6674} & 0.6711 & $\mathbf{-0.0037}$ & 6/15 & 0.89 \\
\bottomrule
\end{tabular}
\end{table}

\subsection{The Protocol Change Is Worth Ten Times More to the Causal Arm}
\label{sec:asym}

The decisive statistic is a different one, and it is decisive because it does not require the two
arms to be comparable. Ask instead how much each arm gains from the protocol change, which is a
within-model difference on identical weights and identical held-out seconds
(Table~\ref{tab:asym}).

\begin{table}[t]
\caption{Protocol Gain by Arm: $\kappa$(Streaming) $-$ $\kappa$(Tiled), Paired Within Fold}
\label{tab:asym}
\centering
\footnotesize
\setlength{\tabcolsep}{5pt}
\begin{tabular}{lcc}
\toprule
 & Sleep-EDF-78 & DOD-H \\
\midrule
Causal arm gains       & $+0.0350$ & $+0.0345$ \\
Non-causal arm gains   & $+0.0036$ & $+0.0052$ \\
\midrule
\textbf{Excess to the causal arm} & $\mathbf{+0.0314}$ & $\mathbf{+0.0292}$ \\
Folds with positive excess & 15/15 & 15/15 \\
Bootstrap 95\% CI & $[.0291, .0334]$ & $[.0253, .0332]$ \\
Corrected $t(14)$ & $12.69$ & $6.48$ \\
$p$ & $4.6\times10^{-9}$ & $1.4\times10^{-5}$ \\
Per-seed range & $.0310$--$.0319$ & $.0285$--$.0297$ \\
\bottomrule
\end{tabular}
\end{table}

The excess is positive in every one of the thirty paired folds across the two datasets, the
smallest single-fold value anywhere being $+0.0173$, and the per-seed spread within a dataset is
under $0.001\,\kappa$. Two corpora that share almost nothing but the label convention agree on the
size of the effect to within $0.002\,\kappa$.

The aggregate cost then follows by arithmetic rather than as an independent observation: a tiled
penalty of $-0.0249$ and $-0.0329$, an excess of $+0.0314$ and $+0.0292$, and a streaming residual
of $+0.0065$ and $-0.0037$. The protocol accounts for essentially all of the measured penalty on
Sleep-EDF-78 and about 89\% of it on DOD-H.

\subsection{The Mechanism, Verified}
\label{sec:latency}

Define permitted latency $L$: the label for second $t$ is emitted at $t + L$, so the model sees $L$
seconds of future and $L_{\text{buf}} - L$ of past, and $L = 0$ is real-time operation. One forward
pass over a window yields a prediction at every position, and position $p$ carries exactly
$L_{\text{buf}} - 1 - p$ seconds of lookahead, so a single sweep measures every latency at once.

Averaging the resulting curve over all 120 buffer positions reproduces that model's tiled score to
within $0.003\,\kappa$ in all four model--dataset combinations: $0.6626$ against a measured
$0.6634$ and $0.6890$ against $0.6912$ on Sleep-EDF-78, $0.6537$ against $0.6567$ and $0.6915$
against $0.6922$ on DOD-H. The tiled protocol is, arithmetically, this average, and that is the
whole explanation.

Table~\ref{tab:quartile} locates the penalty, and the two datasets show the same shape. At buffer
positions carrying 91--120 seconds of history the two arms are level. Across the last quartile,
where a window has under 30 seconds of history, the causal arm loses $0.114$ and $0.128\,\kappa$
against its own best quartile while the non-causal arm loses only $0.018$ and $0.028$, because it
still has up to 119 seconds of lookahead there. That gap is the artifact. A streaming system
occupies those positions once, at start-up, and never again.

\begin{table}[t]
\caption{$\kappa$ by Buffer Quartile, Seed 42, 120~s Buffer. Positive $\Delta$ Means the Causal Arm
Pays}
\label{tab:quartile}
\centering
\footnotesize
\setlength{\tabcolsep}{3.6pt}
\begin{tabular}{lccc|ccc}
\toprule
 & \multicolumn{3}{c|}{Sleep-EDF-78} & \multicolumn{3}{c}{DOD-H} \\
History (s) & Caus. & Non-c. & $\Delta$ & Caus. & Non-c. & $\Delta$ \\
\midrule
91--120 & \textbf{0.6982} & 0.6930 & $-0.0052$ & 0.6904 & 0.6961 & $+0.0057$ \\
61--90  & 0.6923 & 0.6969 & $+0.0046$ & 0.6887 & 0.7044 & $+0.0157$ \\
31--60  & 0.6753 & 0.6911 & $+0.0158$ & 0.6737 & 0.6970 & $+0.0233$ \\
1--30   & 0.5847 & 0.6752 & $+0.0905$ & 0.5622 & 0.6685 & $+0.1063$ \\
\bottomrule
\end{tabular}
\end{table}

Fig.~\ref{fig:latency} gives the full curves. The causal arm declines monotonically from $L = 0$ on
both datasets: it cannot read future signal, so larger $L$ only costs it history, and at $L = 119$,
with one second of history, it collapses to $\kappa = 0.317$ and $0.252$. The non-causal arm traces
an inverted U, gaining lookahead and losing history, and peaks at $L = 46$~s on both datasets, which
is itself a small piece of evidence that the curve shape is a property of the protocol rather than
of either corpus.

\begin{figure}[t]
\centering
\includegraphics[width=\columnwidth]{figures/fig_latency.png}
\caption{Cohen's $\kappa$ against permitted latency $L$, both arms, 5-point moving average over
$L$. Larger $L$ means less history and more lookahead within a fixed 120~s buffer. The causal
collapse at large $L$ is the cold-start behaviour that the tiled protocol averages into its score.}
\label{fig:latency}
\end{figure}

At each arm's best operating point the two datasets disagree, and we report that rather than the
agreement because it bounds the claim. On Sleep-EDF-78 the causal arm reaches $\kappa = 0.7008$ at
zero latency against a non-causal optimum of $0.6977$ at 38~s of delay, a margin of
$0.003\,\kappa$ we would not defend as more than a tie. On DOD-H the causal arm reaches $0.6922$ at
zero latency and the non-causal optimum is $0.7058$ at 46~s, exceeding it by $0.0136$; the
crossover latency is 24~s on Sleep-EDF-78 and 5~s on DOD-H. What replicates across the two datasets is the
protocol asymmetry. The stronger statement, that a non-causal model cannot beat a causal one at any
latency, holds on Sleep-EDF-78 and fails on DOD-H, and a designer should read the first as general
and the second as dataset-specific.

\subsection{The REM Concentration Is the Same Artifact}

Under tiling, REM appears to pay $0.060$ F1 on both datasets, four times the next-largest class on
Sleep-EDF-78 and the only class whose penalty there survives Holm correction across the five
per-class tests. Under the deployment protocol the difference is $+0.002$ and $+0.023$, and N1, N2
and N3 invert on both datasets, the causal arm ending ahead on all three
(Table~\ref{tab:perclass}). This is mechanistically sensible: REM is defined partly by the context
surrounding it, which is why a technician resolves REM against neighbouring epochs rather than
scoring it in isolation, so it is the class most damaged by a nearly empty buffer.

\begin{table}[t]
\caption{Per-Class F1 Difference by Protocol, 15 Paired Folds Per Column. Positive Means the Causal
Arm Pays}
\label{tab:perclass}
\centering
\footnotesize
\setlength{\tabcolsep}{3.2pt}
\begin{tabular}{lcc|cc}
\toprule
 & \multicolumn{2}{c|}{Sleep-EDF-78} & \multicolumn{2}{c}{DOD-H} \\
Class & Tiled & Streaming & Tiled & Streaming \\
\midrule
W   & $+0.018$ & $+0.006$ & $+0.046$ & $+0.035$ \\
N1  & $+0.025$ & $-0.009$ & $+0.007$ & $-0.018$ \\
N2  & $+0.002$ & $-0.017$ & $+0.014$ & $-0.004$ \\
N3  & $+0.008$ & $-0.020$ & $+0.015$ & $-0.004$ \\
REM & $\mathbf{+0.060}$ & $\mathbf{+0.002}$ & $\mathbf{+0.060}$ & $\mathbf{+0.023}$ \\
\midrule
Macro & $+0.023$ & $-0.007$ & $+0.028$ & $+0.006$ \\
\bottomrule
\end{tabular}
\end{table}

Two residuals are worth naming rather than smoothing over. REM still costs the causal arm $0.023$
F1 on DOD-H at deployment, roughly ten times what it costs on Sleep-EDF-78, and Wake costs $0.035$
on DOD-H under both protocols and almost nothing on Sleep-EDF-78. The Wake difference is plausibly
a consequence of the dataset rather than of causality: DOD-H is not Wake-trimmed, so its Wake is
sleep-adjacent and transitional throughout, and transitional Wake is exactly the case where seeing
forward helps. We offer that as a reading, not a measurement.

\subsection{Why the Choice of Estimand Matters at This Sample Size}
\label{sec:estimand}

DOD-H's fold-to-fold $\kappa$ standard deviation is $0.098$ against Sleep-EDF-78's $0.026$. The
same fifteen runs cannot resolve the between-arm causality difference ($p = 0.22$) and do resolve
the within-arm protocol gain ($p = 1.4\times10^{-5}$), because the latter scores identical weights
on identical seconds and so removes model variance from the comparison entirely.

The same lesson appears within Sleep-EDF. On the 78-subject set the tiled causality effect is
$-0.0249\,\kappa$, corrected $t(14) = -3.36$, $p = 0.005$; on the 20-subject split the identical
experiment yields $+0.0143$ with $t(4) = 1.26$, $p = 0.276$, the wrong sign and not significant,
because fold-to-fold variance there is 3.1 times larger. Causality penalties measured on 20-to-25
subject splits, common in this literature for compute reasons, warrant caution: the true effect can
be undetectable or sign-reversed from variance alone.

\subsection{Deployment Behaviour}

Two properties matter for a device and neither is captured by $\kappa$.

\emph{Reporting latency.} A per-second output rate is only useful if it lets a system report
something sooner. Charging both systems for their evidence, a change counts as reported once the
new stage has been held for $h$ seconds, while an epoch-level system cannot emit the label for
epoch $k$ until epoch $k$ has ended, giving it a structural 30-second floor on any change that
begins an epoch. At $h = 10$~s, excluding changes the model was already calling before the search
window opened, the per-second system reports first in 94.8\% of transitions on Sleep-EDF-78 and
95.6\% on DOD-H, with a median 19 seconds saved on both, and on Sleep-EDF-78 the advantage holds
across all five stages entered (+19 to +34 seconds). The result is conditional on $h$ and we report
the full sweep: the share falls to 66.5\% at $h = 20$ and 29.6\% at $h = 30$, where the epoch
system leads, because requiring 30 seconds of confirmation reimposes the epoch and discards the
resolution. A finer rate is worth having precisely when the application will act on less than an
epoch of evidence. Note also that the scored boundary is itself located only to the nearest 30~s,
so these figures describe latency against the scored transition, not the physiological one.

\emph{Stability.} Labelling every second independently makes the raw output far more fragmented
than a scored hypnogram: 167 stage changes an hour against 13 for the technician on Sleep-EDF-78. A
trailing-window mode filter, $\hat{y}_t = \mathrm{mode}(\hat{y}'_{t-w+1},\dots,\hat{y}'_t)$, uses
only past predictions and so preserves causality, needs no retraining, and adds no parameters. At
$w = 30$~s it is worth $+0.0102\,\kappa$ and cuts fragmentation sevenfold to 25 changes an hour; at
$w = 45$~s it reaches 18 an hour, close to the human rate, at negligible cost. N1 F1 peaks earlier,
near $w = 10$~s, since N1 bouts are short and heavy smoothing absorbs them, so the best window
depends on whether overall agreement or transitional-stage sensitivity matters more. The window was
selected on the same held-out folds it is reported on, so this gain is a best-case estimate, and
the sweep was run on Sleep-EDF-78 only.

\emph{Cost.} On one thread of an Intel Core i9-14900HX, batch size 1, median of 1{,}000 timed runs
after 100 warm-up runs, the model costs $0.0271$~ms per second of EEG at 30{,}757 parameters. This
is per-output throughput rather than per-invocation latency; a device recomputing every second pays
a rolling cost roughly forty times larger. MAC counts, peak memory and measurements on
wearable-class silicon are not reported, so this is evidence of feasibility rather than a
device-level result.

\section{Discussion}

An earlier version of this work reported that the causal constraint costs $0.029\,\kappa$ and that
the cost concentrates in REM. Both statements are reproducible under the protocol that produced
them, and both are largely artifacts of it. We report them alongside the corrected measurement
rather than quietly replacing them, because the discrepancy is the finding.

The mechanism is not subtle once seen. A streaming model is defined by what it can do with the
history it has accumulated. Tiling a recording with independent context windows repeatedly destroys
that history and then scores the model on the wreckage, and it does so asymmetrically: a model that
can read forward is largely unaffected at the start of a window, and a model that cannot is
crippled there.

Three things bound the claim. The streaming protocol denies the non-causal arm the lookahead that
defines it, so this compares deployment options rather than architectures; at matched latency the
answer is dataset-dependent, with lookahead worth nothing on Sleep-EDF-78 and about $0.014\,\kappa$
at 46 seconds of delay on DOD-H, and worth roughly $0.025$ to $0.033\,\kappa$ if latency is
unbounded. The non-causal arm is our own matched twin rather than a state-of-the-art offline model,
so the comparison bounds what causality costs within this capacity class. And both datasets are
single-channel EEG scored to AASM conventions with one architecture; the mechanism argues that any
model whose predictions depend on accumulated history should show the effect, but two instances of
one architecture is not a general result.

For a device the practical reading is straightforward. A system that must operate in real time
should not expect to pay for it in accuracy at this capacity scale, should use the streaming
protocol to measure itself, and should be reported against a matched twin rather than against
offline models that differ in a dozen ways at once. For the field the reading is that published
causality penalties should be read together with the evaluation protocol that produced them, and
that reporting both protocols costs nothing, since no retraining is involved.

\section{Conclusion}

A strictly causal per-second sleep-staging model, verified end to end from raw sample to logit, was
compared against a non-causal twin identical in every other respect on two datasets. Under the
batch-tiling protocol standard in this literature the causal constraint costs $0.025$ and
$0.033\,\kappa$ and appears to fall on REM; most of that is produced by the protocol. Moving the
same weights to a streaming protocol lifts the causal arm roughly ten times as much as the
non-causal arm, an excess of $+0.0314$ and $+0.0292\,\kappa$ positive in all 30 paired folds, and
what remains is $+0.0065$ and $-0.0037$, neither resolvable. The tiled score is arithmetically the
mean over context-buffer positions, verified to within $0.003\,\kappa$ in four model--dataset
combinations, and the penalty it reports comes from positions a streaming system occupies once at
start-up. The model runs in 30{,}757 parameters at $0.026$~ms per second of EEG and reports stage
changes a median 19 seconds sooner than an epoch-level system. Code, splits, per-fold predictions
and the analysis for both datasets are released so the measurement can be repeated without
retraining.

\section*{Acknowledgment}

The authors thank Dr.\ Bhaskarjyoti Das for supervision and guidance throughout this project.

\begin{thebibliography}{99}

\bibitem{deepsleepnet} A. Supratak, H. Dong, C. Wu, and Y. Guo, ``DeepSleepNet: A model for
automatic sleep stage scoring based on raw single-channel EEG,'' \emph{IEEE Trans. Neural Syst.
Rehabil. Eng.}, vol. 25, no. 11, pp. 1998--2008, 2017.

\bibitem{attnsleep} E. Eldele, Z. Chen, C. Liu, M. Wu, C.-K. Kwoh, X. Li, and C. Guan, ``An
attention-based deep learning approach for sleep stage classification with single-channel EEG,''
\emph{IEEE Trans. Neural Syst. Rehabil. Eng.}, vol. 29, pp. 809--818, 2021.

\bibitem{usleep} M. Perslev, S. Darkner, L. Kempfner, M. Nikolic, P. J. Jennum, and C. Igel,
``U-Sleep: Resilient high-frequency sleep staging,'' \emph{npj Digit. Med.}, vol. 4, art. 72, 2021.

\bibitem{utime} M. Perslev, M. Jensen, S. Darkner, P. J. Jennum, and C. Igel, ``U-Time: A fully
convolutional network for time series segmentation applied to sleep staging,'' in \emph{Advances in
Neural Information Processing Systems 32}, 2019.

\bibitem{physionet} A. L. Goldberger et al., ``PhysioBank, PhysioToolkit, and PhysioNet: Components
of a new research resource for complex physiologic signals,'' \emph{Circulation}, vol. 101, no. 23,
pp. e215--e220, 2000.

\bibitem{kemp} B. Kemp, A. H. Zwinderman, B. Tuk, H. A. C. Kamphuisen, and J. J. L. Oberye,
``Analysis of a sleep-dependent neuronal feedback loop: the slow-wave microcontinuity of the EEG,''
\emph{IEEE Trans. Biomed. Eng.}, vol. 47, no. 9, pp. 1185--1194, 2000.

\bibitem{dod} A. Guillot, F. Sauvet, E. H. During, and V. Thorey, ``Dreem Open Datasets:
Multi-scored sleep datasets to compare human and automated sleep staging,'' \emph{IEEE Trans.
Neural Syst. Rehabil. Eng.}, vol. 28, no. 9, pp. 1955--1965, 2020.

\bibitem{aasm} R. B. Berry et al., \emph{The AASM Manual for the Scoring of Sleep and Associated
Events: Rules, Terminology and Technical Specifications, Version 2.6}. Darien, IL: American Academy
of Sleep Medicine, 2020.

\bibitem{bresch} E. Bresch, U. Gro\ss{}ekath\"ofer, and G. Garcia-Molina, ``Recurrent deep neural
networks for real-time sleep stage classification from single channel EEG,'' \emph{Front. Comput.
Neurosci.}, vol. 12, art. 85, 2018.

\bibitem{causalattennet} J. Pan, Y. Feng, P. Zhao, X. Zou, A. Hou, and X. Che, ``CausalAttenNet: A
fast and long-term-temporal network for automatic sleep staging with single-channel EEG,''
\emph{IEEE Trans. Instrum. Meas.}, vol. 73, pp. 1--13, 2024.

\bibitem{microsleepnet} G. Liu et al., ``Micro SleepNet: Efficient deep learning model for mobile
terminal real-time sleep staging,'' \emph{Front. Neurosci.}, vol. 17, art. 1218072, 2023.

\bibitem{tcn} S. Bai, J. Z. Kolter, and V. Koltun, ``An empirical evaluation of generic
convolutional and recurrent networks for sequence modeling,'' \emph{arXiv:1803.01271}, 2018.

\bibitem{attention} A. Vaswani et al., ``Attention is all you need,'' in \emph{Advances in Neural
Information Processing Systems 30}, 2017, pp. 5998--6008.

\bibitem{focal} T.-Y. Lin, P. Goyal, R. Girshick, K. He, and P. Doll\'ar, ``Focal loss for dense
object detection,'' in \emph{Proc. IEEE Int. Conf. Comput. Vis. (ICCV)}, 2017, pp. 2980--2988.

\bibitem{nadeau} C. Nadeau and Y. Bengio, ``Inference for the generalization error,'' \emph{Mach.
Learn.}, vol. 52, no. 3, pp. 239--281, 2003.

\end{thebibliography}

\end{document}
